% ============================================================
%  Quatrième de couverture — Quoi que nous soyons
%  Cédric Laubscher — 2026
% ============================================================
\documentclass[12pt,a4paper]{article}
\usepackage[utf8]{inputenc}
\usepackage[T1]{fontenc}
\usepackage[french]{babel}
\usepackage{csquotes}
\usepackage[sfdefault]{sourcesanspro}
\usepackage[top=3.5cm,bottom=3.5cm,left=3cm,right=3cm]{geometry}
\usepackage{amsmath}
\setlength{\parindent}{0pt}
\setlength{\parskip}{0.9em}
\pagestyle{empty}
\begin{document}

\begin{center}
{\large\itshape Quoi que nous soyons}\\[0.3em]
{\small Une introduction au Logo Dynamique Projectif}
\end{center}

\bigskip

Pourquoi un électron existe-t-il~? Non pas comment il se comporte~--- la physique répond à cette question avec une précision stupéfiante. Mais \emph{pourquoi} cette forme-là, cette stabilité absolue, cette identité parfaite à travers le temps et l'espace~? Ce \emph{pourquoi} est la fissure que la physique n'a pas su refermer depuis un siècle.

Le Logo Dynamique Projectif suspend tout~: l'espace, le temps, les particules. Et il pose le problème dans sa nudité absolue~: dans ces conditions, qu'est-ce qui peut exister~? La réponse est simple et vertigineuse. Pour qu'une chose existe durablement, il faut et il suffit qu'elle forme un cycle~--- une différence qui revient sur elle-même. Un battement. De ce seul principe, sans supposer ni l'espace ni le temps, le LDP reconstruit la forme logique de l'électron, puis celle du proton, puis les grandes constantes de la physique~--- $\hbar$, $\alpha$, $k_{\rm B}$, $G$~--- non plus comme des chiffres tombés du ciel, mais comme les signatures de la façon dont chaque structure gère son propre équilibre interne.

Mais aucune structure complexe ne peut se refermer parfaitement sur elle-même. Il reste toujours une fuite incompressible~: une ouverture irréductible vers ce qui la dépasse. C'est de cette fuite que naît la gravitation. C'est cette même ouverture qui rend possibles la liberté, la relation~--- et notre vulnérabilité.

Un électron, un proton, un être humain~: les instruments nécessaires d'une même logique implacable. Ce livre en propose la première lecture unifiée, depuis les axiomes jusqu'à la condition humaine.

\bigskip

\begin{center}
{\small\itshape Toutes les publications techniques du LDP sont disponibles en accès libre sur \textbf{Zenodo}, sous le nom de Cédric Laubscher.}
\end{center}

\end{document}